% --- LaTeX Research Paper Template - S. Venkatraman ---

% --- Set document class and font size ---

\documentclass[letterpaper, 11pt]{article}

% --- Package imports ---

\usepackage{
  amsmath, amsthm, amssymb, mathtools, dsfont, units,   % Math typesetting
  graphicx, wrapfig, subfig, float,                     % Figures and graphics formatting
  listings, color, inconsolata, pythonhighlight,        % Code formatting
  fancyhdr, sectsty, hyperref, enumerate, enumitem }   	% Headers/footers, section fonts, links, lists

% lipsum is just for generating placeholder text and can be removed
\usepackage{lipsum}

% --- Page layout settings ---

% Set page margins
\usepackage[left=1.15in, right=1.15in, top=1.2in, bottom=1in, headsep=.25in]{geometry}

% Anchor footnotes to the bottom of the page
\usepackage[bottom]{footmisc}

% Set line spacing
\renewcommand{\baselinestretch}{1.12}

% Set spacing between paragraphs
\setlength{\parskip}{1.5mm}

% Allow multi-line equations to break onto the next page
\allowdisplaybreaks

% --- Page formatting settings ---

% Set font sizes for section and subsection titles
\sectionfont{\fontsize{13}{15}\selectfont}
\subsectionfont{\fontsize{11}{15}\selectfont}

% Set reference section title font size the same as the section font size
\renewcommand{\refname}{\fontsize{13}{15}\selectfont\bf{References}}

% Set link colors for labeled items and citations (blue) and citations (red)
\hypersetup{colorlinks=true, linkcolor=blue, citecolor=blue}

% Abstract settings
\renewenvironment{abstract}
{\small\begin{quote}\noindent \par{\sc \abstractname.}}
{\noindent\end{quote}}

% --- Settings for printing computer code ---

% Define colors for green text (comments), grey text (line numbers),
% green frame around code, 
\definecolor{greenText}{rgb}{0.5, 0.7, 0.5}
\definecolor{greyText}{rgb}{0.5, 0.5, 0.5}
\definecolor{codeFrame}{rgb}{0.5, 0.7, 0.5}

% Define code settings
\lstdefinestyle{code} {
  frame=single, rulecolor=\color{codeFrame},            % Include a green frame around the code
  numbers=left,                                         % Include line numbers
  numbersep=8pt,                                        % Add space between line numbers and frame
  numberstyle=\tiny\color{greyText},                    % Line number font size (tiny) and color (grey)
  commentstyle=\color{greenText},                       % Put comments in green text
  basicstyle=\linespread{1.1}\ttfamily\footnotesize,    % Set code line spacing
  keywordstyle=\ttfamily\footnotesize,                  % No special formatting for keywords
  showstringspaces=false,                               % No marks for spaces
  xleftmargin=1.95em,                                   % Align code frame with main text
  framexleftmargin=1.6em,                               % Extend frame left margin to include line numbers
  breaklines=true,                                      % Wrap long lines of code
  postbreak=\mbox{\textcolor{greenText}{$\hookrightarrow$}\space} % Mark wrapped lines with an arrow
}

% Set all code listings to be styled with the above settings
\lstset{style=code}

% --- Math/Statistics commands ---

% Add a reference number to a single line of a multi-line equation
% Usage: "\numberthis\label{labelNameHere}" in an align or gather environment
\newcommand\numberthis{\addtocounter{equation}{1}\tag{\theequation}}

% Shortcut for bold text in math mode, e.g. $\b{X}$
\let\b\mathbf

% Shortcut for bold Greek letters, e.g. $\bg{\beta}$
\let\bg\boldsymbol

% Shortcut for calligraphic script, e.g. %\mc{M}$
\let\mc\mathcal

% \mathscr{(letter here)} is sometimes used to denote vector spaces
\usepackage[mathscr]{euscript}

% Convergence: right arrow with optional text on top
% E.g. $\converge[w]$ for weak convergence
\newcommand{\converge}[1][]{\xrightarrow{#1}}

% Normal distribution: arguments are the mean and variance
% E.g. $\normal{\mu}{\sigma}$
\newcommand{\normal}[2]{\mathcal{N}\left(#1,#2\right)}

% Uniform distribution: arguments are the left and right endpoints
% E.g. $\unif{0}{1}$
\newcommand{\unif}[2]{\text{Uniform}(#1,#2)}

% Independent and identically distributed random variables
% E.g. $ X_1,...,X_n \iid \normal{0}{1}$
\newcommand{\iid}{\stackrel{\smash{\text{iid}}}{\sim}}

% Equality: equals sign with optional text on top
% E.g. $X \equals[d] Y$ for equality in distribution
\newcommand{\equals}[1][]{\stackrel{\smash{#1}}{=}}

% Math mode symbols for common sets and spaces. Example usage: $\R$
\newcommand{\R}{\mathbb{R}}   % Real numbers
\newcommand{\C}{\mathbb{C}}   % Complex numbers
\newcommand{\Q}{\mathbb{Q}}   % Rational numbers
\newcommand{\Z}{\mathbb{Z}}   % Integers
\newcommand{\N}{\mathbb{N}}   % Natural numbers
\newcommand{\F}{\mathcal{F}}  % Calligraphic F for a sigma algebra
\newcommand{\El}{\mathcal{L}} % Calligraphic L, e.g. for L^p spaces

% Math mode symbols for probability
\newcommand{\pr}{\mathbb{P}}    % Probability measure
\newcommand{\E}{\mathbb{E}}     % Expectation, e.g. $\E(X)$
\newcommand{\var}{\text{Var}}   % Variance, e.g. $\var(X)$
\newcommand{\cov}{\text{Cov}}   % Covariance, e.g. $\cov(X,Y)$
\newcommand{\corr}{\text{Corr}} % Correlation, e.g. $\corr(X,Y)$
\newcommand{\B}{\mathcal{B}}    % Borel sigma-algebra

% Other miscellaneous symbols
\newcommand{\tth}{\text{th}}	% Non-italicized 'th', e.g. $n^\tth$
\newcommand{\Oh}{\mathcal{O}}	% Big-O notation, e.g. $\O(n)$
\newcommand{\1}{\mathds{1}}		% Indicator function, e.g. $\1_A$

% Additional commands for math mode
\DeclareMathOperator*{\argmax}{argmax}    % Argmax, e.g. $\argmax_{x\in[0,1]} f(x)$
\DeclareMathOperator*{\argmin}{argmin}    % Argmin, e.g. $\argmin_{x\in[0,1]} f(x)$
\DeclareMathOperator*{\spann}{Span}       % Span, e.g. $\spann\{X_1,...,X_n\}$
\DeclareMathOperator*{\bias}{Bias}        % Bias, e.g. $\bias(\hat\theta)$
\DeclareMathOperator*{\ran}{ran}          % Range of an operator, e.g. $\ran(T) 
\DeclareMathOperator*{\dv}{d\!}           % Non-italicized 'with respect to', e.g. $\int f(x) \dv x$
\DeclareMathOperator*{\diag}{diag}        % Diagonal of a matrix, e.g. $\diag(M)$
\DeclareMathOperator*{\trace}{trace}      % Trace of a matrix, e.g. $\trace(M)$

% Numbered theorem, lemma, etc. settings - e.g., a definition, lemma, and theorem appearing in that 
% order in Section 2 will be numbered Definition 2.1, Lemma 2.2, Theorem 2.3. 
% Example usage: \begin{theorem}[Name of theorem] Theorem statement \end{theorem}
\theoremstyle{definition}
\newtheorem{theorem}{Theorem}[section]
\newtheorem{proposition}[theorem]{Proposition}
\newtheorem{lemma}[theorem]{Lemma}
\newtheorem{corollary}[theorem]{Corollary}
\newtheorem{definition}[theorem]{Definition}
\newtheorem{example}[theorem]{Example}
\newtheorem{remark}[theorem]{Remark}

% --- Left/right header text (to appear on every page) ---

% Do not include a line under header or above footer
\pagestyle{fancy}
\renewcommand{\footrulewidth}{0pt}
\renewcommand{\headrulewidth}{0pt}

% Center header text: short paper title
\chead{\MakeUppercase{\scriptsize Singular Value Decomposition Notes}}

% Empty left and right header text
\lhead{}\rhead{}

% --- Document starts here ---

\begin{document}

% --- Title, author, date ---

\title{\normalsize\MakeUppercase{\bfseries 
Singular Value Decomposition Notes}}
% \author{\small\MakeUppercase{
% Your name here}}
\date{\footnotesize\MakeUppercase\today}
\maketitle
\vspace{-1cm}

% --- Main content: import sections as subfiles ---

% \input{../tex_settings/abstract}
% \input{../tex_settings/introduction}

\section{Singular Value Decomposition in Matrix Version}

\begin{lemma}[Rank of $A^T A$ and $A$]
  Let $A$ be a real matrix of size $p \times n$ with rank($A$) = $m$. Then rank($A^T A$) = rank($A$) = $m$.
\end{lemma}
\begin{proof}
  Let $v \in ker(A^T A)$. Then $A^T A v = 0$. Then taking the inner product with $v$, we have $v^T A^T A v = 0$. Then 
  \[
    | Av |^2 = 0 \implies Av = 0 \implies v \in ker(A)
  \] 
  Therefore, $ker(A^T A) \subseteq ker(A)$. Conversely, let $v \in ker(A)$. Then 
  \[
    Av = 0 \implies A^T A v = 0 \implies v \in ker(A^T A)
  \] 
  Therefore, $ker(A) \subseteq ker(A^T A)$. Hence, $ker(A^T A) = ker(A)$ and rank($A^T A$) = rank($A$).
\end{proof}

\begin{lemma}[Spectral Theorem for Symmetric Matrices]\label{lem:spectral_theorem_for_symmetric_matrices}
  Let $P \in M_{n \times n}(\R)$ be a symmetric matrix. Then there exists an orthogonal matrix $Q \in M_{n \times n}(\R)$ such that $Q^T P Q = D$ where $D \in M_{n \times n}(\R)$ is a diagonal matrix with the eigenvalues of $P$ on the diagonal.
\end{lemma}

\begin{theorem}
  Let $A \in M_{p \times n}(\R)$ be a real matrix with rank($A$) = $m$. Then there exists an orthogonal matrix $B \in M_{p \times m}(\R)$ and an orthogonal matrix $C \in M_{n \times m}(\R)$ such that 
  \[
    A = B D C^T,
  \]
  where $D \in M_{m \times m}(\R)$ is a diagonal matrix with positive diagonal entries.
\end{theorem}
\begin{proof}
  We will construct the matrices $B$, $C$, and $D$ explicitly.\\
  Let $(\sigma_1, \sigma_2, \ldots, \sigma_m)$ be the singular values of $A$.
  Define $D$ to be the diagonal matrix with the singular values on the diagonal:
  \[
    D = \text{diag}(\sigma_1, \sigma_2, \ldots, \sigma_m).
  \]
  Since $rank(A) = m$, we know that $\sigma_1, \sigma_2, \ldots, \sigma_m$ are positive singular values of $A$ and thus $D$ is invertible.\\
  Since $A^TA$ is symmetric, there exists an orthogonal matrix $Q \in M_{n \times n}(\R)$ such that $Q^T A^T A Q = E$, then $A^TA = Q E Q^T$ where $E$ is a diagonal matrix with the eigenvalues of $A^TA$ on the diagonal, by Lemma~\ref{lem:spectral_theorem_for_symmetric_matrices}. Define
  \[
    Q = \begin{pmatrix} q_1 & q_2 & \ldots & q_n \end{pmatrix}, \quad \text{where } q_i \text{ are the eigenvectors of } A^TA.
  \]
  Then we can define $C$ be the matrix of the first $m$ columns of $Q$. By definition of $Q$, $C^T C = I_m$.\\
  Next, define $N = D^{-1} = \text{diag}(\sigma_1^{-1}, \sigma_2^{-1}, \ldots, \sigma_m^{-1}) \in M_{m \times m}(\R)$ and $B = A \cdot C \cdot N$.\\
  Now we will show that $B^T B = I_m$.\\
  \begin{align*}
    B^TB &= (A C N)^T (A C N) = N^T C^T A^T A C N \\
         &= N^T C^T Q E Q^T C N = N^T D^2 N = I_m
  \end{align*}
  Note that $C^T Q = \begin{pmatrix} I_m & 0 \end{pmatrix}$ and $Q^T C = \begin{pmatrix} I_m \\ 0 \end{pmatrix}$, so $C^T Q E Q^T C$ picks out the top-left $m \times m$ block of $E$, which equals $D^2 = \text{diag}(\sigma_1^2, \ldots, \sigma_m^2)$. Therefore, $N^T D^2 N = D^{-1} D^2 D^{-1} = I_m$.\\
  Finally, we will show that $A = B D C^T$.
  \[
    B D C^T = (A C N) D C^T = A C (N D) C^T = A C I_m C^T = A C C^T
  \]
  We know that $A^T A = Q E Q^T$ and $C$ is the first $m$ columns of $Q$, then $A^T A = C D^2 C^T$. Then we see
  \[
    A^T A C C^T = C D^2 C^T C C^T = C D^2 I_m C^T = C D^2 C^T = A^T A
  \]
  Thus
  \[
    A^T A (I_n - C C^T) = 0 \implies \text{range}(I - C C^T) \subseteq \ker(A^T A) = \ker(A) \; (\text{by Lemma 1})
  \]
  Thus
  \[
    A (I_n - C C^T) = 0 \implies A - ACC^T = 0 \implies A = ACC^T.
  \]
  Therefore, $A = B D C^T$.
\end{proof}

% --- Bibliography ---

% \bibliography{../tex_settings/References}
% \bibliographystyle{apalike2}

% --- Document ends here ---

\end{document}